\documentclass{article}
\usepackage[utf8]{inputenc}
\usepackage{graphicx}
\usepackage{amsmath}
\usepackage{float} % Add the float package

\title{Mathematical Statistics Assignment 1}
\author{ABID ALI (Student ID: 2023280099)}
\date{October 20, 2024}

\begin{document}

\maketitle

\newpage % Forces the following content to appear on a new page

\section*{1.4.3}

The data in Table 1.9 are presented to illustrate the role of renewable energy consumption in the U.S. energy supply in 2007 (source: \texttt{http://www.eia.doe.gov/fuelrenewable.html}). Renewable energy consists of biomass, geothermal energy, hydroelectric energy, solar energy, and wind energy.

\begin{itemize}
    \item[(a)] Construct a bar graph.
    \item[(b)] Construct a Pareto chart.
    \item[(c)] Construct a pie chart.
\end{itemize}

% Table 1.9 Image placed below the Pareto chart
\begin{figure}[H] % The H option forces the figure to be placed exactly here
\centering
% Replace this with the correct path to your Table 1.9 image file
\includegraphics[width=0.6\textwidth]{D:/Masters NPU Semeter/3rd Semester/Statistics/Homework/Submit/Homework 1/figs/143.png}
\caption{Table 1.9: Renewable Energy Consumption}
\end{figure}

\section*{Solution:}

(a) Bar graph
\begin{figure}[H] % Use the H option to force it to stay here
\centering
% Replace this with the correct path to your bar graph image file
\includegraphics[width=0.8\textwidth]{D:/Masters NPU Semeter/3rd Semester/Statistics/Homework/Submit/Homework 1/figs/bar.png}
\caption{Bar Graph of Energy Consumption}
\end{figure}

(b) Pareto chart
\begin{figure}[H] % Use the H option here as well
\centering
% Replace this with the correct path to your Pareto chart image file
\includegraphics[width=0.8\textwidth]{D:/Masters NPU Semeter/3rd Semester/Statistics/Homework/Submit/Homework 1/figs/pareto.png}
\caption{Pareto Chart}
\end{figure}

(c) Pie chart
\begin{figure}[H] % Use the H option here as well
\centering
% Replace this with the correct path to your pie chart image file
\includegraphics[width=0.8\textwidth]{D:/Masters NPU Semeter/3rd Semester/Statistics/Homework/Submit/Homework 1/figs/pie.png}
\caption{Pie Chart}
\end{figure}

\section*{1.4.14}
The following table gives radon concentrations in pCi/liter (picocurie per liter) obtained from 40 houses in a certain area.

\begin{figure}[H] % The H option forces the figure to be placed exactly here
\centering
% Replace this with the correct path to your Table 1.9 image file
\includegraphics[width=0.8\textwidth]{D:/Masters NPU Semeter/3rd Semester/Statistics/Homework/Submit/Homework 1/figs/1414.png}
\end{figure}

\begin{itemize}
    \item[(a)] Construct a stem-and-leaf display.
    \item[(b)] Construct a frequency histogram and interpret.
    \item[(c)] Construct a pie chart and interpret.
\end{itemize}

\section*{Solution:}
(a) Stem and leaf
\begin{figure}[H] % The H option forces the figure to be placed exactly here
\centering
\begin{minipage}{\textwidth} % Use minipage to ensure the figures stay on the same page
    \centering
    % Replace this with the correct path to your first stem-and-leaf image file
    \includegraphics[width=0.8\textwidth]{D:/Masters NPU Semeter/3rd Semester/Statistics/Homework/Submit/Homework 1/figs/stem_leaf.png}
    \includegraphics[width=0.8\textwidth]{D:/Masters NPU Semeter/3rd Semester/Statistics/Homework/Submit/Homework 1/figs/stem_leaf1.png}
\end{minipage}
\end{figure}
\newpage
(b) Histogram and Interpret
\begin{figure}[H] % Use the H option to force it to stay here
\centering
% Replace this with the correct path to your bar graph image file
\includegraphics[width=0.7\textwidth]{D:/Masters NPU Semeter/3rd Semester/Statistics/Homework/Submit/Homework 1/figs/1414_histogram.png}
\caption{Bar Graph of Energy Consumption}
\end{figure}
(c) Pie chart and Interpret
\begin{figure}[H] % Use the H option to force it to stay here
\centering
% Replace this with the correct path to your bar graph image file
\includegraphics[width=0.7\textwidth]{D:/Masters NPU Semeter/3rd Semester/Statistics/Homework/Submit/Homework 1/figs/1414_pie.png}
\end{figure}

\section*{1.5.5}
Maximal static inspiratory pressure (PI\textsubscript{max}) is an index of respiratory muscle strength. The following data show the measure of PI\textsubscript{max} (cm H\textsubscript{2}O) for 15 cystic fibrosis patients.

\begin{figure}[H] % Use the H option to force it to stay here
\centering
% Replace this with the correct path to your bar graph image file
\includegraphics[width=0.8\textwidth]{D:/Masters NPU Semeter/3rd Semester/Statistics/Homework/Submit/Homework 1/figs/155.png}
\end{figure}

\begin{itemize}
    \item[(a)] Find the lower and upper quartiles, median, and interquartile range. Check for any outliers and interpret.
    \item[(b)] Construct a box plot and interpret.
    \item[(c)] Are there any outliers?
\end{itemize}
\section*{Solution:}

\begin{enumerate}
    \item[(a)] \textbf{Solution:}
    \begin{itemize}
        \item Lower Quartile (Q1) = 80
        \item Median (Q2) = 95
        \item Upper Quartile (Q3) = 115
        \item Interquartile Range (IQR) = Q3 - Q1 = 115 - 80 = 35
        \item Outlier Boundaries: 
        \[
        \text{Lower Bound} = Q1 - 1.5 \times IQR = 80 - 1.5 \times 35 = 27.5
        \]
        \[
        \text{Upper Bound} = Q3 + 1.5 \times IQR = 115 + 1.5 \times 35 = 167.5
        \]
        Since all data points fall between 27.5 and 167.5, \textbf{there are no outliers}.
    \end{itemize}

    \item[(b)] \textbf{Box Plot:}
    \begin{center}
    \includegraphics[width=0.5\textwidth]{D:/Masters NPU Semeter/3rd Semester/Statistics/Homework/Submit/Homework 1/figs/box.png}.
    \end{center}

    \item[(c)] \textbf{Outliers:}
    Based on the lower and upper bounds for outliers (27.5 and 167.5), there are no outliers in this dataset.
\end{enumerate}

\section*{1.5.7}

\noindent (a) For any grouped data with $l$ classes with group frequencies $f_i$ and class midpoints $m_i$, show that
\[
\sum_{i=1}^{l} f_i(m_i - \bar{x}) = 0.
\]

\noindent (b) Verify this result for the data given in \textit{Exercise 1.5.6}.

\begin{figure}[H] % Use the H option to force it to stay here
\centering
% Replace this with the correct path to your bar graph image file
\includegraphics[width=0.7\textwidth]{D:/Masters NPU Semeter/3rd Semester/Statistics/Homework/Submit/Homework 1/figs/156.png}
\caption{Exercise 1.5.6-TABLE 1.21 Class and Frequency.}
\end{figure}

\section*{Solution:}

\textbf{Solution:}

\textbf{(a)} For any grouped data with \( l \) classes and group frequencies \( f_i \) and class midpoints \( m_i \), we need to show that:

\[
\sum_{i=1}^{l} f_i (m_i - \bar{x}) = 0.
\]

The formula for the mean \( \bar{x} \) of grouped data is given by:

\[
\bar{x} = \frac{\sum_{i=1}^{l} f_i m_i}{\sum_{i=1}^{l} f_i}.
\]

We begin by expanding the left-hand side:

\[
\sum_{i=1}^{l} f_i (m_i - \bar{x}) = \sum_{i=1}^{l} f_i m_i - \sum_{i=1}^{l} f_i \bar{x}.
\]

Since \( \bar{x} \) is a constant, we can factor it out from the summation:

\[
\sum_{i=1}^{l} f_i \bar{x} = \bar{x} \sum_{i=1}^{l} f_i.
\]

Now substitute the expression for \( \bar{x} \):

\[
\bar{x} = \frac{\sum_{i=1}^{l} f_i m_i}{\sum_{i=1}^{l} f_i}.
\]

Thus:

\[
\sum_{i=1}^{l} f_i \bar{x} = \frac{\sum_{i=1}^{l} f_i m_i}{\sum_{i=1}^{l} f_i} \cdot \sum_{i=1}^{l} f_i = \sum_{i=1}^{l} f_i m_i.
\]

Now substitute back into the original equation:

\[
\sum_{i=1}^{l} f_i (m_i - \bar{x}) = \sum_{i=1}^{l} f_i m_i - \sum_{i=1}^{l} f_i m_i = 0.
\]

Thus, we have shown that:

\[
\sum_{i=1}^{l} f_i (m_i - \bar{x}) = 0.
\]

\textbf{(b)} Verify the result for the data in Exercise 1.5.6.

The data provided in Exercise 1.5.6 is:

\[
\begin{array}{|c|c|}
\hline
\text{Class} & \text{Frequency} \\
\hline
0-4 & 5 \\
5-9 & 14 \\
10-14 & 15 \\
15-19 & 10 \\
20-24 & 6 \\
\hline
\end{array}
\]

First, calculate the midpoints \( m_i \) for each class:

\[
m_1 = \frac{0+4}{2} = 2, \quad m_2 = \frac{5+9}{2} = 7, \quad m_3 = \frac{10+14}{2} = 12,
\]
\[
m_4 = \frac{15+19}{2} = 17, \quad m_5 = \frac{20+24}{2} = 22.
\]

Next, calculate the mean \( \bar{x} \) of the data:

\[
\bar{x} = \frac{\sum_{i=1}^{5} f_i m_i}{\sum_{i=1}^{5} f_i} = \frac{5(2) + 14(7) + 15(12) + 10(17) + 6(22)}{5 + 14 + 15 + 10 + 6}.
\]

First, compute the numerator:

\[
5(2) + 14(7) + 15(12) + 10(17) + 6(22) = 10 + 98 + 180 + 170 + 132 = 590.
\]

Now, compute the denominator:

\[
5 + 14 + 15 + 10 + 6 = 50.
\]

Thus, the mean \( \bar{x} \) is:

\[
\bar{x} = \frac{590}{50} = 11.8.
\]

Now, verify the result \( \sum_{i=1}^{5} f_i (m_i - \bar{x}) = 0 \) by calculating each term:

\[
f_1(m_1 - \bar{x}) = 5(2 - 11.8) = 5(-9.8) = -49,
\]
\[
f_2(m_2 - \bar{x}) = 14(7 - 11.8) = 14(-4.8) = -67.2,
\]
\[
f_3(m_3 - \bar{x}) = 15(12 - 11.8) = 15(0.2) = 3,
\]
\[
f_4(m_4 - \bar{x}) = 10(17 - 11.8) = 10(5.2) = 52,
\]
\[
f_5(m_5 - \bar{x}) = 6(22 - 11.8) = 6(10.2) = 61.2.
\]

Summing these values:

\[
-49 - 67.2 + 3 + 52 + 61.2 = 0.
\]

Thus, we have verified that:

\[
\sum_{i=1}^{5} f_i (m_i - \bar{x}) = 0.
\]


\section*{1.5.10}
The radon concentration (in pCi/liter) data obtained from 40 houses in a certain area are given below.
\begin{figure}[H] % Use the H option to force it to stay here
\centering
% Replace this with the correct path to your bar graph image file
\includegraphics[width=0.8\textwidth]{D:/Masters NPU Semeter/3rd Semester/Statistics/Homework/Submit/Homework 1/figs/1510.png}
\end{figure}

\begin{enumerate}
    \item[(a)] Find the mean, variance, and range for these data and interpret.
    
    \item[(b)] Find lower and upper quartiles, median, and interquartile range. Check for any outliers.
    
    \item[(c)] Construct a box plot and interpret.
    
    \item[(d)] Construct a histogram.
    
    \item[(e)] Locate on your histogram $\bar{x}$, $\bar{x} \pm s$, $\bar{x} \pm 2s$, and $\bar{x} \pm 3s$. Count the data points in each of the intervals $\bar{x} \pm s$, $\bar{x} \pm 2s$, and $\bar{x} \pm 3s$ and compare this with the empirical rule.
    
\end{enumerate}

\section*{Solution:}

\begin{enumerate}
    \item[(a)] \textbf{Mean, Variance, and Range}
    \begin{itemize}
        \item Mean: $\bar{x} = 7.66$ pCi/L
        \item Variance: $\sigma^2 = 24.39$
        \item Range: $16.7$ pCi/L
    \end{itemize}
    
    \item[(b)] \textbf{Quartiles, Median, and Interquartile Range (IQR)}
    \begin{itemize}
        \item Lower Quartile (Q1): 6.65
        \item Median: 7.6
        \item Upper Quartile (Q3): 11.125
        \item Interquartile Range (IQR): $Q3 - Q1 = 4.475$
        \item Outliers: No outliers were detected.
    \end{itemize}

    \item[(c)] \textbf{Box Plot:}
    \begin{center}
        \includegraphics[width=0.6\textwidth]{D:/Masters NPU Semeter/3rd Semester/Statistics/Homework/Submit/Homework 1/figs/1510_box.png} % Replace with the path to your box plot image
    \end{center}
    
    \item[(d)] \textbf{Histogram:}
    \begin{center}
        \includegraphics[width=0.6\textwidth]{D:/Masters NPU Semeter/3rd Semester/Statistics/Homework/Submit/Homework 1/figs/1510_histogram.png} % Replace with the path to your histogram image
    \end{center}

    \item[(e)] \textbf{Empirical Rule Intervals}
    \begin{itemize}
        \item Mean $\pm$ 1 standard deviation: (2.73, 12.60) — 22 data points
        \item Mean $\pm$ 2 standard deviations: (-2.21, 17.54) — 34 data points
        \item Mean $\pm$ 3 standard deviations: (-7.15, 22.48) — 34 data points
    \end{itemize}
    According to the empirical rule:
    \begin{itemize}
        \item Approximately 68\% (22 data points) lie within 1 standard deviation.
        \item Approximately 95\% (34 data points) lie within 2 standard deviations.
        \item Approximately 99.7\% (34 data points) lie within 3 standard deviations.
    \end{itemize}

\end{enumerate}
\end{document}
